\section{Week 4: caller traces and memory timeline}

\subsection{Data-flow target}

The intended composition is
\[
\begin{array}{c}
  (vk,sk)
  \longrightarrow \mathsf{export}(vku,sku)
  \longrightarrow \mathsf{import\_sk}(sku)
  \longrightarrow \mathsf{hash}_{\rm sign} \\
  \hfill \equiv \mathsf{hash}_{\rm verify}
  \longleftarrow \mathsf{raw\_vk}(vku).
\end{array}
\]
At the already proved local hash boundary this would yield
\[
  \take_{32}(\mu_{\rm sign}^{64})=\mu_{\rm verify}^{32}.
\]

\subsection{KeyGen}

\breakablesame{keypair_raw_api_exact_export_trace} places the actual
\breakablett{cryptolab_haetae_mode2_keypair_internal} on the left of an exact
equivalence and preserves its final memory.  The isolated generated exporter
theorems prove a 992-byte VK write, a 1408-byte SK write, and preservation of
the earlier VK range under the displayed VK/SK disjointness premise.
\breakablesame{keypair_raw_api_exports_matching_prefixes_from_copy_posts}
combines those postconditions with the returned-array prefix theorem.

The attempted single caller Hoare theorem reached the final post-export
continuation but its whole-array existential postcondition caused unbounded
normalization under both \texttt{auto} and \texttt{wp}.  A constructed-memory
replacement was rejected because it would evade the research question.
The exact remaining edge is
\claimid{keypair_raw_api_exports_matching_prefixes_residual}; therefore the
KeyGen caller claim is \status{partial}.

\subsection{Sign}

\theoryname{sign_raw_api_exact_mu_trace} directly relates the actual raw Sign
caller to a trace mirror and preserves final memory and return value.  The
mirror performs the actual signature-core continuation after recording the
imported SK and mu arguments.  The compiled companion theorems prove import
length 1408, usable prefix length 992, and exact propagation of raw API
pre/message addresses and lengths.  This is stronger than a standalone
observer that stops after hashing.

\subsection{Verify and the early-return boundary}

Exact trace equivalences compile for \procname{_sign_verify_tail_m23},
\procname{_verify_full_m23}, \procname{_verify_full_mode2}, and
\procname{sign_verify_internal_mode2_jazz}.  However, the public/raw caller
may reject before the tail: signature unpacking, norm rejection, and the final
mismatch test all precede the mu call.  Hence \(\mathit{siglen}=1474\) is
necessary but insufficient.  The missing lifts are recorded as
\claimid{verify_raw_api_exact_mu_trace_residual_obligation} and
\claimid{cryptolab_verify_internal_exact_mu_trace_residual_obligation}; an
unconditional actual-caller hash-reach theorem would overclaim.

\subsection{Ownership and aliasing}

\begin{center}
\small
\begin{tabular}{p{0.22\linewidth}p{0.31\linewidth}p{0.39\linewidth}}
\toprule
Regions & Classification & Reason \\
\midrule
$vku,sku$ & MUST-DISJOINT & SK export must frame the prior VK export. \\
$seedu$, outputs & MAY-ALIAS-BECAUSE-COPIED-FIRST & Seed is fully imported before output writes. \\
$sku,sigu$ & MAY-ALIAS locally; MUST-DISJOINT for reuse & Sign imports first, but a later Verify experiment needs stable keys. \\
$vku,sigu/siglenu$ & MUST-DISJOINT for composition & Sign output must preserve the VK later read by Verify. \\
$preu,message$ & READ-ONLY-ALIAS-ALLOWED & Both are hash inputs; lengths still define distinct byte streams. \\
$rndu,sigu$ & MAY-ALIAS-BECAUSE-COPIED-FIRST & Randomness is copied before signature export. \\
$sigu,siglenu$ & MUST-DISJOINT & The length write must not corrupt the signature. \\
\bottomrule
\end{tabular}
\end{center}

Cross-call stability is an experiment contract, not a cryptographic
assumption.  EasyCrypt's total memory map is not treated as a memory-safety
proof.
